\chapter{Herramientas y Tecnología de Soporte para SocialScrum}
\label{anexo:herramientas}

Este anexo presenta las adaptaciones contextuales, herramientas tecnológicas y plantillas operacionales recomendadas para implementar \gls{socialscrum} en diversos entornos.

\section{Adaptación contextual y escalabilidad}

SocialScrum debe adaptarse al contexto de cada equipo. El principio general es la \textbf{proporcionalidad}: los equipos pequeños requieren una implementación ligera; los equipos grandes, una estructura más formal.

\subsection{Configuración por tamaño de equipo}

La Tabla~\ref{tab:config-por-tamano} resume la configuración recomendada según el tamaño del equipo.

\begingroup
\scriptsize
\renewcommand{\arraystretch}{1.3}
\setlength{\tabcolsep}{4pt}
\setlength{\LTleft}{0pt}
\setlength{\LTright}{0pt}

\begin{longtable}{|
    >{\RaggedRight\hspace{0pt}}p{3.5cm}|
    >{\RaggedRight\hspace{0pt}}p{3.2cm}|
    >{\RaggedRight\hspace{0pt}}p{3.5cm}|
    >{\RaggedRight\hspace{0pt}}p{4.6cm}|
    }
    \caption{Configuración de SocialScrum por tamaño de equipo.} \label{tab:config-por-tamano}                      \\
    \hline
    \textbf{Aspecto}     & \textbf{Pequeño (5--7)}   & \textbf{Mediano (8--15)}  & \textbf{Grande (16+)}            \\
    \hline \endfirsthead
    \hline
    \multicolumn{4}{|c|}{{\bfseries \tablename\ \thetable{} -- Continuación}}                                         \\
    \hline
    \textbf{Aspecto}     & \textbf{Pequeño (5--7)}   & \textbf{Mediano (8--15)}  & \textbf{Grande (16+)}            \\
    \hline \endhead
    \hline \multicolumn{4}{|r|}{\textit{Continúa...}}                                                               \\ \hline \endfoot
    \hline \endlastfoot
    \gls{social-guide}   & Rol rotativo (10~\%)      & Dedicación 50~\%          & Dedicación 100~\%                \\ \hline
    \gls{health-score}   & 3 dimensiones             & 5 dimensiones             & 5 dim. + agregado organizacional \\ \hline
    Social Retrospective & 15 min integrada          & 30--45 min                & Por subgrupo + general           \\ \hline
    Sobrecarga/Sprint    & 2,0--2,5 h                & 21--29 h                  & 40+ h                            \\ \hline
    Auditoría            & SM informal               & Auditor externo semestral & Auditor externo semestral        \\ \hline
    Escalamiento         & Si HS $<$ 4 por 3 Sprints & Si HS $<$ 4 por 2 Sprints & Siempre dedicado                 \\
\end{longtable}
\endgroup

En \textbf{equipos pequeños}, el Social Guide rotativo permite que todos desarrollen sensibilidad social, distribuyendo la responsabilidad. Requiere capacitación básica de 2~h para todos y un \textit{handover} de 30~min entre Social Guides. En \textbf{equipos medianos}, la dedicación del 50~\% mantiene al Social Guide conectado con la realidad técnica. Para \textbf{equipos grandes}, las opciones son: (a)~un Social Guide por equipo al 50~\%, (b)~un Social Guide compartido al 100~\% para 2--3 equipos, o (c)~un Social Guide Lead con representantes por equipo en organizaciones de 50+ personas.

\subsection{Adaptación por modalidad de trabajo}

El principio para equipos híbridos es \textit{remote-first}: todas las prácticas se diseñan como si todos fueran remotos.

\begingroup
\scriptsize
\renewcommand{\arraystretch}{1.3}
\setlength{\tabcolsep}{4pt}
\setlength{\LTleft}{0pt}
\setlength{\LTright}{0pt}

\begin{longtable}{|
    >{\RaggedRight\hspace{0pt}}p{3.7cm}|
    >{\RaggedRight\hspace{0pt}}p{3.4cm}|
    >{\RaggedRight\hspace{0pt}}p{3.5cm}|
    >{\RaggedRight\hspace{0pt}}p{4.2cm}|
    }
    \caption{Adaptaciones de SocialScrum según modalidad de trabajo.} \label{tab:adaptaciones-modalidad}                              \\
    \hline
    \textbf{Aspecto}     & \textbf{Co-localizado}           & \textbf{Distribuido}                  & \textbf{Híbrido}                \\
    \hline \endfirsthead
    \hline
    \multicolumn{4}{|c|}{{\bfseries \tablename\ \thetable{} -- Continuación}}                                                           \\
    \hline
    \textbf{Aspecto}     & \textbf{Co-localizado}           & \textbf{Distribuido}                  & \textbf{Híbrido}                \\
    \hline \endhead
    \hline \multicolumn{4}{|r|}{\textit{Continúa...}}                                                                                 \\ \hline \endfoot
    \hline \endlastfoot
    Social Daily Standup & Presencial, 15--18 min           & Asíncrono (Slack) o videollamada      & \textit{Remote-first}           \\ \hline
    Social Retrospective & Presencial con \textit{post-its} & Digital (Miro) + videollamada         & Digital para todos              \\ \hline
    Observación          & Directa y contextual             & Patrones en Slack y 1:1 virtuales     & Combinada (señales digitales)   \\ \hline
    Cohesión             & Almuerzos y eventos presenciales & \textit{Coffee chats} virtuales       & Ambas, priorizando inclusividad \\ \hline
    Mediación            & Preferentemente presencial       & Videollamada; presencial si necesario & Según contexto                  \\ \hline
    Health Score         & En la retrospectiva              & Formulario asíncrono + discusión      & Formulario asíncrono para todos \\
\end{longtable}
\endgroup

Para equipos distribuidos, los desafíos principales son la comunicación asíncrona, la ausencia de señales no verbales y el aislamiento~\cite{wilson2006trust, noll2020motivation}. Las adaptaciones incluyen un Social Daily Standup asíncrono mediante bots (Geekbot, Standuply), retrospectivas híbridas (fase asíncrona de 24~h + fase síncrona de 45~min), y observación de patrones en canales de comunicación. Para equipos híbridos, la regla es que todos se conecten desde su dispositivo individual para igualar la experiencia.

\section{Herramientas y tecnología}

\subsection{Integración con herramientas de gestión de proyectos}

La configuración de SocialScrum en herramientas como Jira o Azure DevOps requiere campos personalizados que permitan distinguir ítems sociales de técnicos. La Tabla~\ref{tab:jira-custom-fields-anexo} presenta los campos recomendados; la configuración equivalente aplica para Azure DevOps mediante \textit{Work Item Types} y campos personalizados.

\begingroup
\scriptsize
\renewcommand{\arraystretch}{1.3}
\setlength{\tabcolsep}{4pt}
\setlength{\LTleft}{0pt}
\setlength{\LTright}{0pt}

\begin{longtable}{|
    >{\RaggedRight\hspace{0pt}}p{3.5cm}|
    >{\RaggedRight\hspace{0pt}}p{3.5cm}|
    >{\RaggedRight\hspace{0pt}}p{8.2cm}|
    }
    \caption{Campos personalizados para SocialScrum en herramientas de gestión.} \label{tab:jira-custom-fields-anexo} \\
    \hline
    \textbf{Campo}        & \textbf{Tipo} & \textbf{Descripción}                                                      \\
    \hline \endfirsthead
    \hline
    \multicolumn{3}{|c|}{{\bfseries \tablename\ \thetable{} -- Continuación}}                                           \\
    \hline
    \textbf{Campo}        & \textbf{Tipo} & \textbf{Descripción}                                                      \\
    \hline \endhead
    \hline \endlastfoot
    Item Type             & Select        & TECHNICAL / SOCIAL                                                        \\ \hline
    Social Points (SSP)   & Number        & Escala 1--13, solo para ítems sociales                                    \\ \hline
    Community Smell       & Select        & \textit{Radio Silence}, \textit{Conflict}, \textit{Burnout}, etc.         \\ \hline
    Health Dimension      & Multi-select  & Apertura, Respeto, Coraje, Compromiso, Foco                               \\ \hline
    Social Guide Assigned & User picker   & Responsable del ítem social                                               \\ \hline
    Severity              & Select        & Leve, Moderada, Severa, Crítica                                           \\
\end{longtable}
\endgroup

\subsubsection{Tipos de incidencias y workflow}

Se recomienda iniciar usando el \textit{Issue Type} \texttt{Story} con la etiqueta \texttt{SOCIAL} (filtrable vía JQL: \texttt{labels = "SOCIAL"}). Al madurar el proceso, se puede migrar a un \textit{Issue Type} dedicado ``Social Item'' con un \textit{workflow} propio.

El \textit{workflow} para ítems sociales sigue cinco estados:

\begin{description}[nosep, leftmargin=1.5cm, style=nextline]
    \item[\textit{Backlog}] Ítem identificado, no priorizado.
    \item[\textit{Ready}] Priorizado para el próximo Sprint (transición durante el Sprint Planning).
    \item[\textit{In Progress}] Social Guide trabajando activamente.
    \item[\textit{Monitoring}] Intervención completa, en observación. Pasa a \textit{Done} cuando la dimensión del Health Score mejora $\geq 1$ punto; regresa a \textit{In Progress} si empeora.
    \item[\textit{Done}] Mejoría confirmada en el Health Score.
\end{description}

\noindent Automatizaciones sugeridas: comentario automático al pasar a \textit{Monitoring} (recordar revisión en retrospectiva) y notificación si un ítem permanece en \textit{In Progress} más de 5~días.

\subsection{Plantillas operacionales}

Las siguientes plantillas estandarizan la captura de información en SocialScrum. La Tabla~\ref{tab:resumen-plantillas} resume las seis plantillas disponibles con sus campos clave, y a continuación se detallan las dos más relevantes de modo condensado.

\begingroup
\scriptsize
\renewcommand{\arraystretch}{1.3}
\setlength{\tabcolsep}{4pt}
\setlength{\LTleft}{0pt}
\setlength{\LTright}{0pt}

\begin{longtable}{|
    >{\RaggedRight\hspace{0pt}}p{3.8cm}|
    >{\RaggedRight\hspace{0pt}}p{3cm}|
    >{\RaggedRight\hspace{0pt}}p{8.4cm}|
    }
    \caption{Resumen de plantillas operacionales de SocialScrum.} \label{tab:resumen-plantillas}                                                                                                                                                                            \\
    \hline
    \textbf{Plantilla}                & \textbf{Frecuencia} & \textbf{Secciones principales}                                                                                                                                                                                \\
    \hline \endfirsthead
    \hline \endlastfoot
    User Story Social                 & Por ítem            & Contexto del \textit{smell} (tipo, severidad, detección, impacto en HS), criterios de aceptación, hipótesis de causa raíz (\gls{five-whys}), estrategia de intervención, confidencialidad, estimación en SSP. \\ \hline
    Mediación de conflictos           & Por conflicto       & Partes involucradas, contexto, sesiones 1:1, sesión conjunta, acuerdos firmados, métricas de éxito (encuesta $\geq 6$/10 + mejora en HS).                                                                     \\ \hline
    Health Score Dashboard            & Semanal             & 5 dimensiones con tendencia y \textit{smells} activos, promedio de HS, capacidad estimada, correlación con métricas técnicas (velocidad, defectos).                                                           \\ \hline
    Retrospectiva social              & Por Sprint          & Pre-retro (HS actual, ítems pendientes), análisis de \textit{smells}, nuevos temas, acciones propuestas con SSP, HS esperado para Sprint N+1, notas confidenciales.                                           \\ \hline
    Reporte para \glspl{stakeholders} & Mensual             & Resumen ejecutivo, tendencia de 3~meses, ítems completados, presupuesto social (\% de capacidad), riesgos/mitigaciones, recomendación (continuar/ajustar/escalar).                                            \\ \hline
    Registro de olores                & Continuo            & Identificación (Sprint, evento, detector), definición (manifestación, causa raíz, tipo, severidad), acción tomada (ítem social, SSP, estrategia), resolución y post-mortem.                                   \\
\end{longtable}
\endgroup

\subsubsection{User Story Social (estructura condensada)}

Cada \textit{User Story} social incluye un título con el formato \texttt{[SOCIAL-XXX] Descripción breve} y las siguientes secciones:

\begin{enumerate}[nosep]
    \item \textbf{Contexto del smell:} tipo, severidad (Leve/Moderada/Severa/Crítica), quién y cuándo lo detectó, manifestación observable (sin juicios de valor), personas y roles afectados, dimensión de HS afectada con \textit{score} actual y objetivo.
    \item \textbf{Criterios de aceptación:} mejora de HS en la dimensión objetivo, desaparición del comportamiento observable, satisfacción post-intervención $\geq 7$/10.
    \item \textbf{Hipótesis de causa raíz:} aplicación de los cinco porqués sobre la manifestación.
    \item \textbf{Estrategia de intervención:} tipo (facilitación, mediación, *coaching*, estructural o escalamiento), pasos propuestos y métricas de seguimiento.
    \item \textbf{Confidencialidad:} nivel (público/equipo/confidencial) y restricciones de acceso.
    \item \textbf{Estimación:} desglose en horas (diagnóstico, sesiones 1:1, facilitación grupal, seguimiento) convertido a SSP (1~SSP = 2~horas).
\end{enumerate}

\subsubsection{Mediación de conflictos (estructura condensada)}

La plantilla de mediación es de confidencialidad alta (solo para el Social Guide y las partes implicadas) e incluye:
\begin{enumerate}[nosep]
    \item \textbf{Partes involucradas:} identificación por iniciales/código y número de observadores afectados.
    \item \textbf{Contexto:} origen del conflicto, primera manifestación, eventos agravantes e impacto en la comunicación, colaboración y entregas.
    \item \textbf{Sesiones:} 1:1 con cada parte, sesión conjunta de mediación, documentación de acuerdos firmados y seguimiento programado.
    \item \textbf{\gls{definition-of-done}:} proceso documentado, acuerdos consentidos y comunicados (sin detalles privados), seguimiento programado.
    \item \textbf{Métricas de éxito:} encuesta a ambas partes con un resultado $\geq 6$/10 y mejora del HS en la dimensión afectada.
\end{enumerate}

\subsection{Alertas automatizadas}

La Tabla~\ref{tab:alertas} presenta las alertas recomendadas para monitorear la salud social del equipo de forma proactiva.

\begingroup
\scriptsize
\renewcommand{\arraystretch}{1.3}
\setlength{\tabcolsep}{4pt}
\setlength{\LTleft}{0pt}
\setlength{\LTright}{0pt}

\begin{longtable}{|
    >{\RaggedRight\hspace{0pt}}p{4cm}|
    >{\RaggedRight\hspace{0pt}}p{6.9cm}|
    >{\RaggedRight\hspace{0pt}}p{4.3cm}|
    }
    \caption{Alertas automatizadas para SocialScrum.} \label{tab:alertas}                                     \\
    \hline
    \textbf{Alerta}           & \textbf{Condición}                           & \textbf{Canal de notificación} \\
    \hline \endfirsthead
    \hline \endlastfoot
    HS Crítico                & HS Promedio $< 4{,}0$                        & Slack + Email a CTO            \\ \hline
    Dimensión en caída        & Cualquier dim. baja $> 2$ puntos en 1 Sprint & Slack \#social-guide           \\ \hline
    Smell crítico sin asignar & Severity = Crítica en Backlog $> 48$~h       & Email a Social Guide           \\ \hline
    Capacidad social excedida & SSP en Sprint $> 20~\%$ de capacidad total   & Slack \#scrum-master           \\
\end{longtable}
\endgroup

\subsection{Consideraciones de implementación}

Las herramientas no son un requisito estricto para implementar SocialScrum, pero reducen la fricción operacional. El principio rector es que la tecnología debe servir al proceso: si una simple hoja de cálculo funciona para el equipo, no hay necesidad de una infraestructura compleja. La sofisticación debe ser proporcional al valor que genera. Para la visualización del Health Score, cualquier herramienta de \textit{dashboards} (Google Sheets, Grafana, Power BI) puede conectarse a la API de la herramienta de gestión y al formulario del Health Score para generar paneles de tendencia, correlación con la velocidad y estado de los \gls{community-smells}.